\documentclass[10pt]{scrartcl}
\usepackage[mdthm,colorsec,noauthor,noasy]{evan}
\usepackage{handout}
\newcommand{\ip}[2]{\langle #1,#2\rangle}
\DeclareMathOperator{\tr}{tr}\DeclareMathOperator{\nul}{null}\DeclareMathOperator{\range}{range}\DeclareMathOperator{\spn}{span}
\newcommand{\LL}{\mathcal L}\newcommand{\PP}{\mathcal P}\newcommand{\MM}{\mathcal M}
\title{Math 110 --- Week 12}
\subtitle{Solutions}
\author{}\date{}\ihead{}
\begin{document}
\maketitle

\begin{prob}
  \textbf{\xlink{math110}{2026-11-16-positive-isometries-svd}{prob-HW12-P1}{HW12 P1}. }\par
  \idea{For a real symmetric $2\times2$ matrix, $\det$ is the product and $\tr$ is the sum of its two real eigenvalues (\xlink{math110}{2026-11-16-positive-isometries-svd}{prop-105}{Props.~105--106}).}
  \begin{steps}
    \ii $A$: $\det A=3>0$ and $\tr A=4>0$, so both eigenvalues ($3,1$) are strictly positive: \textbf{positive definite}.
    \ii $B$: $\det B=1-4=-3<0$, so the eigenvalues ($3,-1$) have opposite signs: \textbf{not positive semidefinite} ($\langle B(1,-1),(1,-1)\rangle=-2$).
    \ii $C$: $\det C=0$ and $\tr C=2$, so the eigenvalues are $2,0$: \textbf{positive semidefinite, not positive definite} ($\langle C(1,-1),(1,-1)\rangle=0$ with $(1,-1)\neq0$).
  \end{steps}
\end{prob}

\begin{prob}
  \textbf{\xlink{math110}{2026-11-16-positive-isometries-svd}{prob-HW12-P2}{HW12 P2}. }\par
  \idea{Use the $2\times2$ square-root formula or diagonalize orthogonally with $\frac{1}{\sqrt{2}}(1,\pm1)$.}
  Using the \xlink{math110}{2026-11-16-positive-isometries-svd}{name-square-root-2x2}{square-root formula} with $\det A=3$, $\tr A=4$, and $4+2\sqrt{3}=(1+\sqrt{3})^2$:
  \[
    \sqrt{A} = \frac{1}{1+\sqrt{3}}\begin{bmatrix}2+\sqrt{3}&1\\1&2+\sqrt{3}\end{bmatrix}
    = \frac{\sqrt{3}-1}{2}\begin{bmatrix}2+\sqrt{3}&1\\1&2+\sqrt{3}\end{bmatrix}
    = \frac{1}{2}\begin{bmatrix}\sqrt{3}+1&\sqrt{3}-1\\\sqrt{3}-1&\sqrt{3}+1\end{bmatrix}.
  \]
  (Check via orthogonal diagonalization: $A$ has eigenvalues $3,1$ on $\frac{1}{\sqrt{2}}(1,\pm1)$, so $\sqrt{A}$ has eigenvalues $\sqrt{3},1$ on the same vectors, whose square recovers $A$.)
\end{prob}

\begin{prob}
  \textbf{\xlink{math110}{2026-11-16-positive-isometries-svd}{prob-HW12-P3}{HW12 P3}. }\par
  \idea{An isometry preserves inner products; conversely, Pythagoras computes $\|Sv\|$ from orthonormal coordinates.}
  \begin{steps}
    \ii ($\Rightarrow$) If $S$ is an isometry, $S^*S=I$ (\xlink{math110}{2026-11-16-positive-isometries-svd}{prop-109}{Prop.~109}), so $\langle Se_i,Se_j\rangle=\langle S^*Se_i,e_j\rangle=\langle e_i,e_j\rangle=\delta_{ij}$: $Se_1,\dots,Se_n$ is orthonormal.
    \ii ($\Leftarrow$) If $Se_1,\dots,Se_n$ is orthonormal in $W$, write $v=\sum_{k=1}^n \lambda_k e_k$. Then $Sv=\sum_{k=1}^n \lambda_k Se_k$, and Pythagoras / (T2) gives $\|Sv\|^2=\sum_{k=1}^n |\lambda_k|^2 = \|v\|^2$.
  \end{steps}
\end{prob}

\begin{prob}
  \textbf{\xlink{math110}{2026-11-16-positive-isometries-svd}{prob-HW12-P4}{HW12 P4}. }\par
  \idea{$A^{\mathsf T}A$ is diagonal so $P=\sqrt{A^{\mathsf T}A}$ is immediate, and invertibility gives $S=AP^{-1}$.}
  \begin{steps}
    \ii $A^{\mathsf T}A = \operatorname{diag}(1,4)$, so its unique positive square root is $P=\sqrt{A^{\mathsf T}A}=\operatorname{diag}(1,2)$ (\xlink{math110}{2026-11-16-positive-isometries-svd}{prop-107}{Prop.~107}).
    \ii Since $A$ is invertible,
    \[
      S = AP^{-1} = \begin{bmatrix}0&2\\1&0\end{bmatrix}\begin{bmatrix}1&0\\0&1/2\end{bmatrix} = \begin{bmatrix}0&1\\1&0\end{bmatrix},
    \]
    which is orthogonal (reflection across $y=x$). Thus
    \[ A = \begin{bmatrix}0&1\\1&0\end{bmatrix}\begin{bmatrix}1&0\\0&2\end{bmatrix}. \]
  \end{steps}
\end{prob}

\begin{prob}
  \textbf{\xlink{math110}{2026-11-16-positive-isometries-svd}{prob-HW12-P5}{HW12 P5}. }\par
  \idea{Orthonormal eigenvectors of $A^{\mathsf T}A$ give $e_1,e_2$; set $f_1=Ae_1/s_1$ and complete $\{f_1\}$ to an orthonormal basis.}
  \begin{steps}
    \ii $A^{\mathsf T}A=\begin{psmallmatrix}2&2\\2&2\end{psmallmatrix}$ has eigenvalues $4$ on $e_1=\frac{1}{\sqrt{2}}(1,1)$ and $0$ on $e_2=\frac{1}{\sqrt{2}}(1,-1)$, so $s_1=2$, $s_2=0$.
    \ii $f_1 = Ae_1/s_1$, and $Ae_1=\frac1{\sqrt2}(2,2)$ with $s_1=2$, so $f_1=\frac{1}{\sqrt{2}}(1,1)$.
    \ii Extending $\{f_1\}$ to an orthonormal basis gives $f_2=\frac{1}{\sqrt{2}}(1,-1)$, and $Ae_2=0=0\cdot f_2$.
  \end{steps}
\end{prob}

\begin{prob}
  \textbf{\xlink{math110}{2026-11-16-positive-isometries-svd}{prob-HW12-P6}{HW12 P6}. }\par
  \idea{Invert $Te_k=s_kf_k$ and invoke the converse SVD proposition (\xlink{math110}{2026-11-16-positive-isometries-svd}{prop-116}{Prop.~116}).}
  By \xlink{math110}{2026-11-16-positive-isometries-svd}{prop-113}{Prop.~113}, take orthonormal bases $e_1,\dots,e_n$ and $f_1,\dots,f_n$ with $Te_k=s_k f_k$. Since $T$ is invertible, every $s_k=\|Te_k\|>0$. Applying $T^{-1}$ and dividing by $s_k$ gives $T^{-1}f_k = s_k^{-1} e_k$. Since $\{f_k\}$ and $\{e_k\}$ are orthonormal bases and $s_k^{-1}>0$, \xlink{math110}{2026-11-16-positive-isometries-svd}{prop-116}{Prop.~116} identifies the singular values of $T^{-1}$ as $s_1^{-1},\dots,s_n^{-1}$.
\end{prob}

\begin{prob}
  \textbf{\xlink{math110}{2026-11-16-positive-isometries-svd}{prob-HW12-P7}{HW12 P7}. }\par
  \idea{Bound the sum of squares of coefficients in an SVD basis.}
  In an SVD $Te_j=s_jf_j$ (\xlink{math110}{2026-11-16-positive-isometries-svd}{prop-113}{Prop.~113}), write $v=\sum_j c_j e_j$. Then $Tv=\sum_j s_j c_j f_j$, and since $s_j\le s_1$,
  \[
    \|Tv\|^2 = \sum_{j=1}^n s_j^2 |c_j|^2 \le s_1^2 \sum_{j=1}^n |c_j|^2 = s_1^2 \|v\|^2.
  \]
  Taking square roots gives $\|Tv\|\le s_1\|v\|$, with equality at $v=e_1\neq0$ ($\|Te_1\|=\|s_1f_1\|=s_1=s_1\|e_1\|$).
\end{prob}

\begin{prob}
  \textbf{\xlink{math110}{2026-11-16-positive-isometries-svd}{prob-O1}{O1}. }\par
  \idea{$\sqrt{S}$ acts as a scalar on each eigenspace of $S$, which is $T$-invariant because $ST=TS$.}
  \begin{steps}
    \ii $ST$ is self-adjoint: $(ST)^*=T^*S^*=TS=ST$.
    \ii For any $v\in E(\lambda,S)$, $S(Tv)=T(Sv)=\lambda(Tv)$, so $E(\lambda,S)$ is $T$-invariant.
    \ii Since $\sqrt{S}|_{E(\lambda,S)}=\sqrt{\lambda}I$ (\xlink{math110}{2026-11-16-positive-isometries-svd}{prop-107}{Prop.~107}), on $v\in E(\lambda,S)$ we have $\sqrt{S}Tv=\sqrt{\lambda}Tv=T(\sqrt{\lambda}v)=T\sqrt{S}v$. Because $V$ has an eigenbasis of $S$, $\sqrt{S}T=T\sqrt{S}$ everywhere.
    \ii Using self-adjointness of $\sqrt{S}$ and positiveness of $T$,
    \[
      \langle STv,v\rangle = \langle\sqrt{S}\,T\sqrt{S}\,v,v\rangle = \langle T(\sqrt{S}v),\sqrt{S}v\rangle \ge 0.
    \]
    \ii Without $ST=TS$: $S=\operatorname{diag}(1,0)$ and $T=\begin{psmallmatrix}1&1\\1&1\end{psmallmatrix}$ are positive, yet $ST=\begin{psmallmatrix}1&1\\0&0\end{psmallmatrix}$ is not self-adjoint, hence not positive.
  \end{steps}
\end{prob}

\begin{prob}
  \textbf{\xlink{math110}{2026-11-16-positive-isometries-svd}{prob-O2}{O2}. }\par
  \idea{$AA^{\mathsf T}$ is only $2\times2$, so find $f_1,f_2$ first and pull back to $e_j=A^{\mathsf T}f_j/s_j$.}
  \begin{steps}
    \ii $AA^{\mathsf T}=\begin{psmallmatrix}2&1\\1&2\end{psmallmatrix}$ has eigenvalues $3$ on $f_1=\frac{1}{\sqrt{2}}(1,1)$ and $1$ on $f_2=\frac{1}{\sqrt{2}}(1,-1)$. Thus the singular values of $A$ are $\sqrt{3},1,0$.
    \ii Set $e_1 = A^{\mathsf T}f_1/\sqrt{3} = \frac{1}{\sqrt{6}}(1,1,2)$ and $e_2 = A^{\mathsf T}f_2/1 = \frac{1}{\sqrt{2}}(1,-1,0)$, completed to $\RR^3$ by $e_3=\frac{1}{\sqrt{3}}(1,1,-1)\in\nul A$.
    \ii Check $Ae_1=\sqrt{3}f_1$, $Ae_2=f_2$, $Ae_3=0$, so $A=F\Sigma E^{\mathsf T}$ with
    \[
      F = \frac{1}{\sqrt{2}}\begin{bmatrix}1&1\\1&-1\end{bmatrix},\qquad
      \Sigma = \begin{bmatrix}\sqrt{3}&0&0\\0&1&0\end{bmatrix},\qquad
      E = [e_1\;e_2\;e_3].
    \]
  \end{steps}
\end{prob}

\begin{prob}
  \textbf{\xlink{math110}{2026-11-16-positive-isometries-svd}{prob-O3}{O3}. }\par
  \idea{$A^{\mathsf T}A$ has perfect-square determinant $225=15^2$, making both $P=\sqrt{A^{\mathsf T}A}$ and the SVD exact.}
  \begin{steps}
    \ii $A^{\mathsf T}A=\begin{psmallmatrix}25&20\\20&25\end{psmallmatrix}$ has $\tr=50$, $\det=225=15^2$, so by \xlink{math110}{2026-11-16-positive-isometries-svd}{name-square-root-2x2}{Square-root formula for $2\times2$ positive matrices}:
    \[
      P = \frac{A^{\mathsf T}A+15I}{\sqrt{50+30}} = \frac{1}{4\sqrt{5}}\begin{bmatrix}40&20\\20&40\end{bmatrix} = \sqrt{5}\begin{bmatrix}2&1\\1&2\end{bmatrix}, \qquad
      S = AP^{-1} = \frac{1}{3\sqrt{5}}\begin{bmatrix}3&0\\4&5\end{bmatrix}\begin{bmatrix}2&-1\\-1&2\end{bmatrix} = \frac{1}{\sqrt{5}}\begin{bmatrix}2&-1\\1&2\end{bmatrix}.
    \]
    \ii Eigenvalues of $A^{\mathsf T}A$ are $45$ on $e_1=\frac{1}{\sqrt{2}}(1,1)$ and $5$ on $e_2=\frac{1}{\sqrt{2}}(1,-1)$, giving $s_1=3\sqrt{5}$, $s_2=\sqrt{5}$, and
    \[
      f_1 = \frac{Ae_1}{3\sqrt{5}} = \frac{1}{\sqrt{10}}(1,3),\qquad
      f_2 = \frac{Ae_2}{\sqrt{5}} = \frac{1}{\sqrt{10}}(3,-1),
    \]
    yielding $A=F\Sigma E^{\mathsf T}$ with $F=[f_1\;f_2]$, $\Sigma=\operatorname{diag}(3\sqrt{5},\sqrt{5})$, $E=[e_1\;e_2]$.
  \end{steps}
\end{prob}

\begin{prob}
  \textbf{\xlink{math110}{2026-11-16-positive-isometries-svd}{prob-O4}{O4}. }\par
  \idea{A positive operator's singular values equal its eigenvalues.}
  Because $T^*T$ is positive, $\sqrt{(T^*T)^*(T^*T)}=\sqrt{(T^*T)^2}=T^*T$ (\xlink{math110}{2026-11-16-positive-isometries-svd}{prop-107}{Prop.~107}), so the singular values of $T^*T$ are its eigenvalues $s_1^2\ge\dots\ge s_n^2$ (\xlink{math110}{2026-11-16-positive-isometries-svd}{def-54}{Def.~54}). Hence $\|T^*T\|_2 = s_1^2 = \|T\|_2^2$.
\end{prob}

\end{document}
